\section{Introduction}\label{sec:introduction}

Airports are canonical cyber-physical systems: aircraft, ground vehicles, passengers, baggage, and weather interact continuously under tight schedules and regulatory constraints. Digital twins---virtual counterparts that track a physical system through structured data flows~\cite{glaessgen2012,grieves2017}---are increasingly proposed as the operating paradigm for such infrastructure, yet in practice they are frequently delivered as monolithic, closed-loop models coupling simulation logic, geometry, live data feeds, and analytics in a single artefact. The difficulty is structural rather than incidental. Fidelity expectations grow over a project's lifetime: a team starts with logical operational modelling because it is cheap and testable, but stakeholders then demand physical realism, geospatial visualisation, live flight-tracking, and finally decision support. Each concern is individually well understood---network-level delay propagation has been modelled analytically~\cite{pyrgiotis2013}, and the cost of delay minutes has been quantified empirically~\cite{cook2009}---yet accreted onto a monolithic core, every step up in fidelity forces an architectural rewrite. Extensibility and reproducibility remain unresolved challenges~\cite{jones2020}: such a twin is hard to extend and effectively impossible to reproduce, its behaviour depending on live feeds, hidden local state, and undocumented assumptions.

We define digital twin in architectural terms: a graph-persisted representation of a system's entities and their relationships, kept synchronised with the source system's state through a fact-only event stream. Queries against the graph reflect current and historical state with bounded latency and support both monitoring and command issuance back into the source system. KART's ``source system'' is itself a synthetic simulation; the claim is that the pattern generalises to real physical systems. Read against the digital-model/shadow/twin taxonomy~\cite{kritzinger2018}, layers L0--L2 form a digital shadow, and only the decisional layer's approval-gated autonomous commands close the loop into a true digital twin.

The resulting design question is how a digital twin of a complex cyber-physical system should be architected so that fidelity can be acquired progressively without architectural rewrites, and whether specification-first development with AI coding agents can make such a system tractable for a single developer. Answering it requires three things: the layered architecture itself, the specification-and-skill methodology, and a complete end-to-end implementation whose costs, benefits, and failure modes can be examined honestly.

To answer these questions, we designed a four-layer fidelity architecture and instantiated it as Arthur International Airport (KART), a fully synthetic airport digital twin comprising roughly 45{,}000 lines of Python across nine microservices and roughly 28{,}000 lines of React/TypeScript dashboards. The system was developed by a single author over more than fifty sprints using a specification-first methodology: event-bus schemas, a graph data model, simulation rules, and one authoritative specification per service are written before implementation, complemented by ``skill files'' that encode reusable patterns and gotchas for both humans and LLM agents.

This paper makes four contributions:
\begin{enumerate}
  \item a layered fidelity architecture for digital twins in which each capability layer strictly consumes the contracts of the layer below, enabling progressive enrichment without re-platforming;
  \item a specification-first development methodology with skill files tailored to AI coding agents;
  \item KART, a fully open-source airport digital twin spanning flight operations, passenger flow, baggage handling, weather impacts, incident cascades, and financial planning;
  \item an honest account of lessons learned from fifty-plus sprints of agent-assisted development, including defects that slipped through review and the process changes that remediated them.
\end{enumerate}

The remainder of this paper is organised as follows. Section~\ref{sec:related} surveys related work on digital twins, airport simulation, and AI-assisted software engineering. Section~\ref{sec:architecture} presents the layered fidelity architecture. Section~\ref{sec:methodology} describes the specification-first development methodology and its integration with AI coding agents. Section~\ref{sec:implementation} details the KART implementation. Section~\ref{sec:evaluation} reports the evaluation approach and results, Section~\ref{sec:discussion} discusses findings and threats to validity, and Section~\ref{sec:conclusion} concludes.
